\section{Ingineria Interfețelor Mobile și React Native}

Dezvoltarea aplicațiilor mobile moderne impune provocări legate de diversitatea dispozitivelor și a sistemelor de operare. StyleGenAI utilizează framework-ul \textbf{React Native} pentru a oferi o experiență nativă, menținând în același timp o bază de cod unificată.

\subsection{Paradigma Arhitecturală Bridge}
Spre deosebire de aplicațiile hibride clasice (bazate pe WebView), React Native nu randează HTML/CSS. În schimb, acesta invocă componentele native ale platformei gazdă (Android/iOS).
Arhitectura se bazează pe trei fire de execuție paralele:
\begin{itemize}
    \item \textbf{Main Thread (UI)}: Gestionează randarea și interacțiunea directă cu ecranul.
    \item \textbf{JavaScript Thread}: Execută logica de business a aplicației (codul React).
    \item \textbf{Shadow Thread}: Calculează layout-ul elementelor folosind algoritmul \textit{Yoga} (o implementare a Flexbox în C++).
\end{itemize}
Comunicarea între firul JS și cel nativ se realizează asincron printr-un "Bridge" (punte), care serializează mesajele în format JSON. Această decuplare asigură că interfața rămâne fluidă (60 FPS) chiar dacă logica JS este intensivă.

\subsection{Programarea Declarativă și Virtual DOM}
Abordarea React este fundamental declarativă: dezvoltatorul descrie \textit{starea} dorită a interfeței, iar framework-ul calculează tranzițiile necesare.
Conceptul central este \textbf{Virtual DOM}, o reprezentare în memorie a structurii UI.
Algoritmul de \textit{Reconciliation} (Diffing) compară noul arbore Virtual DOM cu cel anterior la fiecare modificare de stare, identificând setul minim de operații native necesare pentru actualizarea ecranului. Aceasta optimizează performanța, evitând redesenarea completă a interfeței.

\subsection{Managementul Stării și Fluxul Unidirecțional}
Stabilitatea aplicației este asigurată de principiul fluxului de date unidirecțional (\textit{One-Way Data Flow}). Datele curg de la componentele părinte spre copii prin proprietăți (\texttt{props}), asigurând predictibilitatea stării.
Pentru starea globală (ex: sesiunea de autentificare, tema vizuală), se utilizează \textbf{Context API}. Acesta implementează pattern-ul \textit{Dependency Injection}, permițând injectarea datelor necesare direct în componentele din adâncimea arborelui, fără a polua componentele intermediare.

\subsection{Integrarea cu Hardware-ul Dispozitivului}
Prin intermediul bibliotecilor \textbf{Expo SDK}, aplicația abstractizează accesul la resursele hardware:
\begin{itemize}
    \item \textbf{Senzori Optici}: Modulul \textit{Expo Camera} oferă control programatic asupra parametrilor de captură (focus, flash, aspect ratio).
    \item \textbf{File System}: Gestionează cache-ul local și manipularea fișierelor binare înainte de upload.
    \item \textbf{Secure Enclave}: Utilizează \textit{Android Keystore} și \textit{iOS Keychain} pentru stocarea criptată a token-urilor JWT, protejând sesiunea utilizatorului chiar și pe dispozitive cu acces root.
\end{itemize}